% !TEX program = xelatex
\documentclass[11pt,a4paper]{article}
\usepackage{algo-preamble}

\title{\textbf{L13 \textperiodcentered{} Άπληστοι Αλγόριθμοι III}\\[2pt]
\large Κωδικοποίηση Huffman}
\author{Algorithms Class Hub}
\date{\small Αλγόριθμοι και Πολυπλοκότητα (K17) \textperiodcentered{} ΕΚΠΑ}

\begin{document}
\maketitle

\begin{center}\itshape\small
Συμπίεση δεδομένων με κωδικοποίηση Huffman: απροθεματικοί κώδικες, αναπαράσταση
με δυαδικά δέντρα, ο άπληστος αλγόριθμος συγχώνευσης συχνοτήτων, και η απόδειξη
βελτιστότητας.
\end{center}

\section{Τι θα δούμε}

Ένα μόνο πρόβλημα --- αλλά από τα ωραιότερα του μαθήματος: \textbf{πώς συμπιέζεις
δεδομένα στο ελάχιστο δυνατό μήκος;} Η απάντηση, η \textbf{κωδικοποίηση Huffman},
είναι ένας καθαρός άπληστος αλγόριθμος με μια κομψή απόδειξη βελτιστότητας μέσω
\textbf{επιχειρήματος ανταλλαγής} (όπως στο προηγούμενο μάθημα).

\section{Το πρόβλημα}

\textbf{Είσοδος:} μια αλφαριθμητική ακολουθία δεδομένων --- π.χ.\ το κείμενο ενός
email.

\textbf{Στόχος:} αναπαράστησε τα δεδομένα με τα δυφία (bits) $0, 1$ στο
\textbf{συντομότερο δυνατό μήκος}.

\begin{warningbox}
\textbf{Δεν είναι κρυπτογράφηση.} Η Huffman είναι \textbf{τεχνική συμπίεσης} ---
στόχος της είναι να μικρύνει το αρχείο, όχι να το κρύψει. Η αποκωδικοποίηση
πρέπει να είναι όχι μόνο εφικτή, αλλά και \textbf{μονοσήμαντη}.
\end{warningbox}

Συγκεκριμένα: έχουμε $n$ διαφορετικούς χαρακτήρες, και για καθέναν ξέρουμε τη
\textbf{συχνότητα} εμφάνισής του $f_1, \dots, f_n$. Θέλουμε να αναθέσουμε σε κάθε
χαρακτήρα μια \textbf{συμβολοσειρά δυφίων}, ώστε το συνολικό μήκος του
κωδικοποιημένου κειμένου να είναι ελάχιστο.

\section{Σταθερό vs μεταβλητό μήκος}

\textbf{Σταθερό μήκος:} δίνουμε σε \textbf{κάθε} χαρακτήρα τον ίδιο αριθμό
δυφίων. Με $k$ δυφία κωδικοποιούνται $2^k$ χαρακτήρες.
\begin{itemize}
  \item $\checkmark$ Εύκολη αποκωδικοποίηση --- κόβεις ανά $k$ δυφία.
  \item $\times$ Δεν πετυχαίνει ελάχιστο μήκος.
\end{itemize}

\textbf{Μεταβλητό μήκος:} δίνουμε \textbf{λίγα} δυφία στους συχνούς χαρακτήρες
και \textbf{περισσότερα} στους σπάνιους.
\begin{itemize}
  \item $\checkmark$ Πετυχαίνει ελάχιστο μήκος.
  \item $\times$ Πιο σύνθετη κωδικοποίηση/αποκωδικοποίηση.
\end{itemize}

\begin{intuition}
\textbf{Το παράδειγμα.} Αρχείο 100.000 χαρακτήρων με συχνότητες (σε χιλιάδες):

\begin{center}
\begin{tabular}{@{}l c c c c c c@{}}
\toprule
Χαρακτήρας       & a & b & c & d & e & f \\
\midrule
Συχνότητα        & 45 & 13 & 12 & 16 & 9 & 5 \\
Σταθερό μήκος    & 000 & 001 & 010 & 011 & 100 & 101 \\
Μεταβλητό μήκος  & 0 & 101 & 100 & 111 & 1101 & 1100 \\
\bottomrule
\end{tabular}
\end{center}

\textbf{Σταθερό μήκος:} $100.000 \times 3 = 300.000$ δυφία. \textbf{Μεταβλητό
μήκος:} $45 \cdot 1 + 13 \cdot 3 + 12 \cdot 3 + 16 \cdot 3 + 9 \cdot 4 + 5 \cdot
4 = 224.000$ δυφία. Κέρδος \textbf{25\%} --- απλώς δίνοντας στο πανταχού παρόν
\texttt{a} έναν κώδικα ενός δυφίου.
\end{intuition}

\section{Το πρόβλημα της αποκωδικοποίησης}

Το μεταβλητό μήκος κρύβει μια παγίδα. Έστω αλφάβητο $a, b, c$ με κώδικες
$a \to 01$, $b \to 010$, $c \to 1$. Σε ποια λέξη αντιστοιχεί η συμβολοσειρά
\texttt{0101};
\begin{itemize}
  \item $aa = 01 \sqcup 01$
  \item $bc = 010 \sqcup 1$
\end{itemize}
\textbf{Διφορούμενο!} Δύο τρόποι, η ίδια συμβολοσειρά. Το πρόβλημα: ο κώδικας του
$a$ (\texttt{01}) είναι \textbf{πρόθεμα} του κώδικα του $b$ (\texttt{010}).

\begin{keypoint}
\textbf{Απροθεματική κωδικοποίηση} (prefix-free code): μια συνάρτηση $c$ που
αντιστοιχεί κάθε χαρακτήρα σε μια συμβολοσειρά δυφίων, έτσι ώστε για κάθε δύο
διαφορετικούς χαρακτήρες $x \ne y$, η $c(x)$ να \textbf{μην} είναι πρόθεμα της
$c(y)$.
\end{keypoint}

Με απροθεματικό κώδικα η αποκωδικοποίηση γίνεται \textbf{μονοσήμαντη}. Π.χ.\ με
$c(a) = 11,\ c(b) = 01,\ c(d) = 001,\ c(e) = 10,\ c(f) = 000$, η συμβολοσειρά
\texttt{11010001010001} διαβάζεται μοναδικά ως \texttt{abfeed}.

\section{Κόστος μιας κωδικοποίησης}

\begin{keypoint}
Το \textbf{κόστος} μιας κωδικοποίησης $c$ είναι ο συνολικός αριθμός δυφίων:
\[
\text{cost}(c) = \sum_{x \in S} f_x \cdot |c(x)|
\]
όπου $|c(x)|$ το μήκος του κώδικα του $x$. Ζητάμε \textbf{απροθεματική}
κωδικοποίηση που \textbf{ελαχιστοποιεί} αυτό το άθροισμα.
\end{keypoint}

\section{Κωδικοποιήσεις ως δυαδικά δέντρα}

Κάθε απροθεματική κωδικοποίηση αντιστοιχεί σε ένα \textbf{δυαδικό δέντρο}:
\begin{itemize}
  \item Οι \textbf{χαρακτήρες} κάθονται στα \textbf{φύλλα}.
  \item Το μονοπάτι από τη ρίζα στο φύλλο δίνει τον κώδικα: \textbf{\texttt{0}
  $=$ αριστερό παιδί, \texttt{1} $=$ δεξιό παιδί}.
\end{itemize}

\begin{intuition}
\textbf{Γιατί αυτό εγγυάται απροθεματικότητα;} Αν ένας χαρακτήρας είναι σε
\textbf{φύλλο}, κανένας άλλος χαρακτήρας δεν βρίσκεται «κάτω» από αυτόν. Άρα
κανένας κώδικας δεν είναι πρόθεμα κάποιου άλλου --- εξ ορισμού του φύλλου.
\textbf{Αποκωδικοποίηση:} διάβαζε δυφία ακολουθώντας το μονοπάτι από τη ρίζα·
μόλις φτάσεις σε φύλλο, βρήκες έναν χαρακτήρα --- ξεκίνα ξανά από τη ρίζα.
\end{intuition}

Σε αυτή τη γλώσσα, το $|c(x)|$ είναι ακριβώς το \textbf{βάθος} $d_T(\ell_x)$ του
φύλλου του $x$, και θέλουμε δέντρο που ελαχιστοποιεί το
$\sum_x f_x \cdot d_T(\ell_x)$.

\section{Η άπληστη ιδέα}

\begin{keypoint}
\textbf{Άπληστο κριτήριο.} Χαρακτήρες με \textbf{μικρές} συχνότητες πάνε σε
φύλλα \textbf{μακριά} από τη ρίζα (μεγάλος κώδικας)· χαρακτήρες με
\textbf{μεγάλες} συχνότητες πάνε σε φύλλα \textbf{κοντά} στη ρίζα (μικρός
κώδικας).
\end{keypoint}

Πώς το χτίζουμε; \textbf{Από κάτω προς τα πάνω.} Παίρνουμε τους δύο
\textbf{σπανιότερους} χαρακτήρες και τους κάνουμε αδέλφια --- αυτοί θα είναι οι
πιο βαθιά. Τους «συγχωνεύουμε» σε έναν εικονικό χαρακτήρα με συχνότητα το
άθροισμά τους, και επαναλαμβάνουμε.

\begin{codebox}
\begin{verbatim}
Huffman(S):                          // S = χαρακτήρες με συχνότητες
  Εάν |S| = 2:
    επίστρεψε δέντρο με ρίζα και δύο φύλλα
  Αλλιώς:
    έστω y, z οι δύο χαρακτήρες με τη μικρότερη συχνότητα στο S
    S' <- S
    αφαίρεσε τα y, z από το S'
    εισήγαγε νέο χαρακτήρα ω στο S' με f_ω = f_y + f_z
    T' <- Huffman(S')
    T <- το T' με την ω να αποκτά δύο παιδιά, τα φύλλα y και z
    επίστρεψε T
\end{verbatim}
\end{codebox}

\begin{intuition}
\textbf{Διαισθητικά:} οι δύο σπανιότεροι χαρακτήρες «αξίζουν» τις χειρότερες
θέσεις. Τους κλειδώνουμε αμέσως ως αδέλφια στο βαθύτερο σημείο, τους
αντιμετωπίζουμε από εκεί και πέρα σαν \textbf{ένα} σύνθετο σύμβολο, και λύνουμε
το ίδιο πρόβλημα με έναν χαρακτήρα λιγότερο.
\end{intuition}

\subsection{Παράδειγμα}

Χαρακτήρες με συχνότητες $a{:}45,\ e{:}12,\ f{:}13,\ g{:}16,\ b{:}5,\ d{:}9$.
Συγχωνεύσεις, κάθε φορά οι δύο μικρότερες: $b + d = 14$· $e + f = 25$·
$14 + g = 30$· $25 + 30 = 55$· $a + 55 = 100$. Το δέντρο που προκύπτει δίνει τους
κώδικες $a = $ \texttt{0}, $e = $ \texttt{100}, $f = $ \texttt{101},
$g = $ \texttt{111}, $b = $ \texttt{1100}, $d = $ \texttt{1101}: συχνά κοντά στη
ρίζα, σπάνια βαθιά.

\subsection{Πολυπλοκότητα}

Ο αλγόριθμος κάνει $n - 1$ συγχωνεύσεις. Κάθε συγχώνευση χρειάζεται να βρει τις
δύο μικρότερες συχνότητες --- με μια \textbf{ουρά προτεραιότητας} (από το μάθημα
των δομών δεδομένων) κοστίζει $O(\log n)$. Άρα
\[
T(n) = T(n - 1) + O(\log n) \ \Longrightarrow\ T(n) = O(n \log n).
\]

\section{Ορθότητα}

Η απόδειξη στηρίζεται σε ένα \textbf{επιχείρημα ανταλλαγής} φύλλων.

\subsection{Λήμμα 1 --- βάθος αντιστρόφως ανάλογο της συχνότητας}

\begin{keypoint}
\textbf{Λήμμα.} Σε ένα \textbf{βέλτιστο} δέντρο $T$, αν δύο χαρακτήρες έχουν
συχνότητες $f(k) \le f(j)$, τότε $d_T(\ell_k) \ge d_T(\ell_j)$ --- ο σπανιότερος
είναι \textbf{τουλάχιστον τόσο βαθιά}.
\end{keypoint}

\textbf{Απόδειξη.} Έστω, για άτοπο, $f(k) < f(j)$ αλλά
$d_T(\ell_k) < d_T(\ell_j)$. \textbf{Ανταλλάσσουμε} τα φύλλα των $k$ και $j$,
παίρνοντας κωδικοποίηση $C'$. Η διαφορά κόστους:
\[
\begin{aligned}
L(C) - L(C') &= f(k)\,d_T(\ell_k) + f(j)\,d_T(\ell_j) - f(k)\,d_T(\ell_j) -
f(j)\,d_T(\ell_k) \\
&= \bigl(f(k) - f(j)\bigr)\bigl(d_T(\ell_k) - d_T(\ell_j)\bigr) \ >\ 0
\end{aligned}
\]
(γινόμενο δύο \textbf{αρνητικών} ποσοτήτων). Άρα $L(C') < L(C)$ --- το $T$
\textbf{δεν ήταν} βέλτιστο. \textbf{Άτοπο.} $\blacksquare$

\subsection{Λήμμα 2 --- δομή του βέλτιστου δέντρου}

\begin{keypoint}
Σε κάθε βέλτιστο δέντρο, κάθε \textbf{εσωτερική} κορυφή έχει \textbf{ακριβώς 2
παιδιά}. Επιπλέον, οι δύο \textbf{μεγαλύτερου μήκους} κώδικες ανήκουν στους
\textbf{δύο σπανιότερους} χαρακτήρες και διαφέρουν \textbf{μόνο στο τελευταίο
δυφίο} (είναι αδέλφια).
\end{keypoint}

Αν μια κορυφή $v$ είχε \textbf{ένα μόνο} παιδί $u$, θα μπορούσαμε να
«συνθλίψουμε» την ακμή $\{v, u\}$ --- μικραίνοντας κάθε κώδικα από κάτω και άρα
το κόστος. Άρα στο βέλτιστο δέντρο αυτό δεν συμβαίνει. Από το Λήμμα 1, οι δύο
σπανιότεροι χαρακτήρες είναι οι βαθύτεροι --- και μπορούμε να τους θεωρήσουμε
\textbf{αδέλφια}.

\subsection{Θεώρημα --- η Huffman είναι βέλτιστη}

\begin{keypoint}
\textbf{Θεώρημα.} Για κάθε πλήθος χαρακτήρων $n$ και κάθε συνάρτηση συχνοτήτων,
το δέντρο που παράγει ο αλγόριθμος Huffman αποτελεί \textbf{βέλτιστη}
κωδικοποίηση.
\end{keypoint}

\textbf{Απόδειξη (επαγωγή στο $|S|$).}
\begin{itemize}
  \item \textit{Βάση} $n = 2$: το δέντρο με μία ρίζα και δύο φύλλα είναι πάντα
  βέλτιστο.
  \item \textit{Επαγωγική υπόθεση:} το θεώρημα ισχύει για $n - 1$ χαρακτήρες.
  \item \textit{Επαγωγικό βήμα:} έστω, για άτοπο, ότι το δέντρο $T$ του
  αλγορίθμου για $n$ χαρακτήρες \textbf{δεν} είναι βέλτιστο. Τότε υπάρχει δέντρο
  $Z$ με $\text{cost}(Z) < \text{cost}(T)$ --- και (Λήμμα 2) μπορούμε να το
  πάρουμε με τους δύο σπανιότερους χαρακτήρες $x, y$ \textbf{αδέλφια}.
\end{itemize}
Από το $Z$ φτιάχνουμε το $Z'$: αφαιρούμε τα $x, y$ και βάζουμε στη θέση του
κοινού γονέα τους έναν χαρακτήρα $\omega$ με $f_\omega = f_x + f_y$. Όμοια, από
το $T$ φτιάχνουμε το $T'$. Κρίσιμη ταυτότητα:
\[
\text{cost}(Z') = \text{cost}(Z) - f_\omega \qquad \text{cost}(T') =
\text{cost}(T) - f_\omega
\]
(όταν «βγάζεις» τα δύο αδέλφια, χάνεις ακριβώς $f_x + f_y$ από το κόστος --- κάθε
ένα έχανε ένα επίπεδο βάθους). Αφού $\text{cost}(Z) < \text{cost}(T)$,
αφαιρώντας το ίδιο $f_\omega$: $\text{cost}(Z') < \text{cost}(T')$.

Αλλά το $T'$ είναι ακριβώς το δέντρο που παράγει ο αλγόριθμος για $n - 1$
χαρακτήρες --- και από την \textbf{επαγωγική υπόθεση} είναι βέλτιστο. Δεν μπορεί
να υπάρχει $Z'$ φθηνότερο. \textbf{Άτοπο.} $\blacksquare$

\begin{intuition}
\textbf{Η ψυχή της απόδειξης:} η άπληστη πρώτη κίνηση --- «κάνε τους δύο
σπανιότερους αδέλφια» --- δεν θυσιάζει τίποτα, γιατί κάποιο βέλτιστο δέντρο
\textbf{ούτως ή άλλως} τους έχει αδέλφια (Λήμμα 2). Άρα μπορούμε με ασφάλεια να
τους κλειδώσουμε και να αναχθούμε στο ίδιο πρόβλημα με έναν χαρακτήρα λιγότερο.
\end{intuition}

\begin{recapbox}
\textbf{Τι κρατάμε από αυτή τη διάλεξη}
\begin{enumerate}
  \item \textbf{Κωδικοποίηση Huffman} --- συμπίεση δεδομένων: λίγα δυφία στους
  συχνούς χαρακτήρες, πολλά στους σπάνιους.
  \item \textbf{Απροθεματικός κώδικας} --- κανένας κώδικας δεν είναι πρόθεμα
  κάποιου άλλου· εγγυάται μονοσήμαντη αποκωδικοποίηση.
  \item Κάθε απροθεματικός κώδικας $=$ \textbf{δυαδικό δέντρο} με χαρακτήρες στα
  φύλλα· κόστος $\sum_x f_x \cdot d_T(\ell_x)$.
  \item \textbf{Άπληστος αλγόριθμος:} συγχώνευσε επανειλημμένα τους \textbf{δύο
  σπανιότερους} χαρακτήρες. $O(n \log n)$ με ουρά προτεραιότητας.
  \item \textbf{Ορθότητα} μέσω επιχειρήματος ανταλλαγής φύλλων $+$ επαγωγής: η
  άπληστη πρώτη κίνηση συμφωνεί με κάποιο βέλτιστο δέντρο, και η ταυτότητα
  $\text{cost}(T') = \text{cost}(T) - f_\omega$ ανάγει το πρόβλημα σε μικρότερο.
\end{enumerate}
\end{recapbox}

\lecturefooter
\end{document}
